%% notes-tex/database/database-expansion-200/sections/summary.tex
%% [[experiments.database.expansion-200]]
%% https://github.com/Mjvolk3/torchcell/tree/main/notes-tex/database/database-expansion-200/sections/summary.tex
%% SOURCE: experiments/database/scripts/build_candidate_datasets_table.py

\clearpage
\section{Summary statistics\secstatus{todo}}\label{sec:summary}

%% GENERATED FILE -- do not hand-edit.
%% SOURCE: experiments/database/scripts/build_candidate_datasets_table.py
\begin{table}[H]\centering
\small
\caption[]{Summary of the final sixty, split at the fifty line. \emph{Top 50} and
\emph{Extra} count rows; \emph{Genotypes}, \emph{Instances} and \emph{Meas.} sum the row
axes over all sixty and are lower bounds, since a row with no count contributes nothing.
The sixty rows name 89 joins, 55 of them to a dataset already built.}
\label{tab:summary}
\begin{tabular}{@{}l r r r r r@{}}
\toprule
 & Top 50 & Extra & Genotypes & Instances & Meas. \\
\midrule
\multicolumn{6}{@{}l}{\emph{By band}}\\
state across genomes & 8 & 0 & 3,000 & $3.1\times 10^{3}$ & $1.2\times 10^{7}$ \\
perturb-seq & 23 & 0 & 327,436 & $4.0\times 10^{5}$ & $4.1\times 10^{6}$ \\
molecular layers & 14 & 0 & 5,500 & $5.7\times 10^{3}$ & $1.1\times 10^{6}$ \\
scale & 5 & 10 & 100,252,284 & $1.5\times 10^{8}$ & $1.5\times 10^{8}$ \\
\midrule
\multicolumn{6}{@{}l}{\emph{By tier}}\\
tier 1 & 8 & 9 & 100,253,896 & $1.1\times 10^{8}$ & $1.2\times 10^{8}$ \\
tier 2 & 18 & 1 & 287,527 & $3.8\times 10^{7}$ & $4.2\times 10^{7}$ \\
tier 3 & 21 & 0 & 46,794 & $4.1\times 10^{4}$ & $5.4\times 10^{6}$ \\
tier 4 & 3 & 0 & 3 & 17 & $5.6\times 10^{4}$ \\
\midrule
\multicolumn{6}{@{}l}{\emph{By class}}\\
CRISPR library screen & 8 & 0 & 307,294 & $3.7\times 10^{5}$ & $3.7\times 10^{5}$ \\
Deep mutational scan & 0 & 2 & 88,000 & $8.8\times 10^{4}$ & $8.8\times 10^{4}$ \\
Expression / single cell & 15 & 0 & 522 & 715 & $4.7\times 10^{6}$ \\
Metabolite / precursor & 5 & 2 & 16,101 & $3.8\times 10^{4}$ & $1.4\times 10^{5}$ \\
Modality / backbone & 4 & 1 & 14,447 & $1.7\times 10^{4}$ & $5.8\times 10^{4}$ \\
Natural variation & 14 & 1 & 120,150 & $1.7\times 10^{6}$ & $1.4\times 10^{7}$ \\
Regulatory DNA & 1 & 0 & 100,000,000 & $1.0\times 10^{8}$ & $1.0\times 10^{8}$ \\
Tolerance / robustness & 3 & 4 & 41,706 & $5.1\times 10^{7}$ & $5.1\times 10^{7}$ \\
\midrule
\multicolumn{6}{@{}l}{\emph{By sequence basis}}\\
S288C+ORF-plasmid & 0 & 1 & 5,280 & $1.1\times 10^{4}$ & $1.1\times 10^{4}$ \\
S288C+designed-edit & 2 & 2 & 164,000 & $1.6\times 10^{5}$ & $1.6\times 10^{5}$ \\
S288C+guide & 7 & 0 & 242,194 & $3.1\times 10^{5}$ & $3.1\times 10^{5}$ \\
S288C+reporter-locus & 1 & 0 & 100,000,000 & $1.0\times 10^{8}$ & $1.0\times 10^{8}$ \\
S288C+tag & 1 & 1 & 8,190 & $2.8\times 10^{4}$ & $3.3\times 10^{4}$ \\
S288C-KO & 8 & 5 & 48,189 & $5.1\times 10^{7}$ & $5.4\times 10^{7}$ \\
engineered-chassis & 1 & 0 & 200 & 200 & 200 \\
isolate-WGS & 5 & 0 & 16,645 & $8.9\times 10^{5}$ & $2.9\times 10^{6}$ \\
reference-only & 16 & 0 & 17 & 243 & $1.1\times 10^{6}$ \\
segregant-WGS & 9 & 1 & 103,505 & $8.2\times 10^{5}$ & $1.1\times 10^{7}$ \\
\midrule
\multicolumn{6}{@{}l}{\emph{Perturb-seq axis}}\\
high on both axes & 1 & 0 & 1,012 & $1.0\times 10^{3}$ & $5.8\times 10^{6}$ \\
input axis only & 5 & 2 & 100,289,094 & $1.0\times 10^{8}$ & $1.0\times 10^{8}$ \\
output axis only & 9 & 0 & 2,779 & $5.4\times 10^{3}$ & $2.4\times 10^{6}$ \\
neither axis & 35 & 8 & 295,335 & $5.2\times 10^{7}$ & $6.0\times 10^{7}$ \\
\midrule
\multicolumn{6}{@{}l}{\emph{Other attributes}}\\
carries a time axis & 6 & 0 & 51 & 77 & $10.0\times 10^{5}$ \\
has a join to a built dataset & 38 & 2 & 100,455,061 & $1.5\times 10^{8}$ & $1.6\times 10^{8}$ \\
has a join to another candidate only & 11 & 2 & 87,236 & $9.2\times 10^{4}$ & $7.3\times 10^{6}$ \\
no named join & 1 & 6 & 45,923 & $1.0\times 10^{5}$ & $1.4\times 10^{6}$ \\
figures sourced this pass & 26 & 8 & 355,000 & $5.2\times 10^{7}$ & $5.4\times 10^{7}$ \\
figures from recall, to confirm & 24 & 2 & 100,233,220 & $1.0\times 10^{8}$ & $1.2\times 10^{8}$ \\
blocked on data access & 1 & 0 & 11,000 & $1.3\times 10^{7}$ & $1.3\times 10^{7}$ \\
loader in flight & 0 & 1 & 5,000 & $3.8\times 10^{7}$ & $3.8\times 10^{7}$ \\
instances a reported count & 26 & 0 & 100,080,698 & $1.1\times 10^{8}$ & $1.3\times 10^{8}$ \\
instances a product of the axes & 6 & 7 & 151,303 & $3.9\times 10^{7}$ & $4.1\times 10^{7}$ \\
instances an estimate & 18 & 3 & 356,219 & $4.4\times 10^{5}$ & $6.5\times 10^{5}$ \\
\midrule
Total & 50 & 10 & 100,588,220 & $1.5\times 10^{8}$ & $1.7\times 10^{8}$ \\
\bottomrule
\end{tabular}
\end{table}


The fifty are the three priority bands in full, 8 state rows, 23 perturb-seq rows and 14
molecular-layers rows, plus the first five scale rows. The ten extra are all scale, and
the fifty's five scale rows include the two largest datasets in the whole ranking, de Boer
2020 and Lee 2014, which is why the instance sums in Table~\ref{tab:summary} are carried
mostly by five rows. The banded rows are vector-valued: a few hundred strains each carrying
a transcriptome, a proteome or a per-cell distribution, so they hold few instances and many
measurements, and ranking on instances alone would have ordered the two blocks the other
way round. The state band is the extreme case, 3,000 genotypes and $1.2\times10^7$
measurements.

Two class facts follow from the bands. Natural variation and expression are the joint
largest classes in the sixty at fifteen rows each, and 16 of the sixty are
\term{reference-only}, meaning the environment rather than a strain edit carries the
perturbation, which is what a condition-crossed expression layer looks like. Thirteen of
the sixty sit on the S288C deletion collection, the axis most of the served set already
uses, so those thirteen are the cheapest joins in the table.

\subsection{Time\secstatus{todo}}\label{sec:time}

Six of the sixty rows carry a time dimension (Table~\ref{tab:final}, \emph{Time} column):
Jariani 2020's lag-phase course, Jackson 2023's continuous course, Wang 2022's three ages,
Airoldi 2016's upshift course, and the two rate rows, Sun 2013 and Martin-Perez 2017,
whose rates are derived from a labeling series that is not itself released. Hackett 2020's
induction series, the row the previous pass's time discussion was anchored at, is served
as 1,752 induction time-course records, one environment per time point. The steady-state
chemostat rows (Brauer 2008, Boer 2010, Hackett 2016, Yu 2021) vary dilution rate, which
is a growth rate, not a time axis, and are not counted.

What the schema holds is \texttt{Environment.duration\_hours} and
\texttt{Environment.duration\_generations}, treatment durations that make two exposure
lengths two distinct environments, which is how Hackett 2020 was ingested. It is a
flattening: a model sees time as a categorical condition, and nothing on the record says
that eight environments are one series in order. A first-class time field on the record
and a series identity that groups its points are the two schema additions a $c(t)$ model
needs, and the six rows above are the ones that would exercise them. They are recorded
here as an open schema item rather than folded into this list's ranking, since the ranking
rule does not yet score a time axis.
